%%%%%%%%%%%%%%%%%%%%%%%%%%%%%%%%%%%%%%%%%%%%%%%%%%%%%%%%%%%
%                           CAPA
%%%%%%%%%%%%%%%%%%%%%%%%%%%%%%%%%%%%%%%%%%%%%%%%%%%%%%%%%%%

%%%%%%%%%%%%%%%%%%%%%%%%% ALTERAR %%%%%%%%%%%%%%%%%%%%%%%%%
\newcommand{\universidade}{Universidade Federal de Minas Gerais}
\newcommand{\faculdade}{Programa de Pós-Graduação em Engenharia de Estrutuas}
\newcommand{\depto}{}
\newcommand{\curso}{}
\newcommand{\disciplina}{EES847 - Análise não-linear geométrico de estruturas reticuladas}
\newcommand{\autor}{Pedro Henrique Reis da Silva Lima}
\newcommand{\titulo}{Trabalhos Práticos da Disciplina de Análise Não-Linear Geométrico de
Estruturas Reticuladas}
                     % | colocar ": " iniciando o subtítulo, caso houver
                     % v
\newcommand{\subtitulo}{}
\newcommand{\local}{Belo Horizonte}
\newcommand{\ano}{\the\year{}}
%%%%%%%%%%%%%%%%%%%%%%% NÃO ALTERAR %%%%%%%%%%%%%%%%%%%%%%%
\hspace{-1.5cm}
\begin{minipage}[t][][c]{\textwidth}
    % Logotipo da Universidade - 
    % - largura da imagem - largura da caixa - comentar para remover
    \includeunivlogo{arquivos/brasao_ufmg.eps}{2.5cm}{2.5cm}
    %    
    \centering{
    \begin{minipage}{10cm}
        \centering
        \textbf{\MakeUppercase{\universidade}}\par      % Comentar para remover
        \textbf{\MakeUppercase{\faculdade}}\par         % Comentar para remover
        %\textbf{\MakeUppercase{\depto}}\par             % Comentar para remover
        %\textbf{\MakeUppercase{\curso}}\par             % Comentar para remover
        \textbf{\MakeUppercase{\disciplina}}            % Comentar para remover  
    \end{minipage}}
    %
    % Logitipo da Unidade - largura da imagem - largura da caixa - comentar para remover
    \includeunitlogo{arquivos/logo_propees.png}{2.5cm}{2.5cm}
    % Segundo a Escola de Engenharia, o Brasão permanece mas deve ser usado em documentos 
    % oficiais e solenes.
\end{minipage}
{
\thispagestyle{empty}
\centering\par
%
\vspace{4em}
%
\autor\par

\vspace{15em}
%%%%%%%%%%%%%%%%%%%%%%%%%%%%%%%%%%%%%%%%%%%%%%%%%%%%%%%
\textbf{\large \MakeUppercase{\titulo}\linebreak\subtitulo}\par
\vfill
\local\par
\ano\par
}